%! TeX root: ../informe.tex
\section{Métodos}
\label{sec:metodos}
En esta sección se detalla la metodología seguida para evaluar tanto el
rendimiento como la calidad de los sistemas de extracción de características y
algoritmos de agrupamiento para nuestro banco de imágenes de grafiti.

\subsection{Arquitectura del entorno de evaluación}
\label{subsec:arquitectura_pipeline}
% Anotación: Describir la arquitectura modular del pipeline implementado.
% Mencionar cómo las clases abstractas en abstractions.py permiten el desacoplamiento.
% Explicar el flujo de datos desde la imagen en bruto hasta el resultado del clustering.

\subsubsection{Diseño general del sistema}
% Anotación: Detallar el uso de Python y herramientas como uv para la gestión
% de dependencias.
% Resaltar la importancia de la reproducibilidad y la automatización mediante
% el pipeline de benchmark.

\subsubsection{Descripción del conjunto de datos}
% Anotación: Describir el dataset de imágenes de grafiti (volumen, tipos de
% imágenes, resolución).
% Detallar los pasos de preprocesamiento aplicados (redimensionamiento,
% normalización, etc.).

\subsection{Espacio de configuraciones del sistema}
\label{subsec:configuraciones}
% Anotación: Introducir que el sistema permite combinar múltiples componentes.

\subsubsection{Modelos de extracción de características (Embeddings)}
% Anotación: Listar y justificar los modelos CNN y YOLO utilizados.
% Mencionar la arquitectura de salida de cada modelo y la dimensionalidad de
% los vectores generados.

\subsubsection{Algoritmos de agrupamiento (Clustering)}
% Anotación: Explicar la selección de algoritmos (K-Means, DBSCAN, HDBSCAN,
% OPTICS).
% Discutir brevemente los hiperparámetros clave que se variarán en las pruebas.

\subsubsection{Técnicas de reducción de dimensionalidad}
% Anotación: Justificar el uso de PCA y UMAP.
% Explicar cómo afectan a la eficiencia del clustering y a la visualización de
% resultados.

\subsubsection{Sistemas de almacenamiento y recuperación}
% Anotación: Comparar el uso de SQLite para pruebas rápidas frente a PostgreSQL
% con pgvector para producción/escalabilidad.

\subsection{Métricas de evaluación}
\label{subsec:metricas}
% Anotación: Definir qué se entiende por "éxito" en este estudio, tanto en
% calidad como en eficiencia.

\subsubsection{Métricas de calidad y desempeño}
% Anotación: Detallar métricas intrínsecas (Silueta, Davies-Bouldin).
% Explicar cómo se evalúa la precisión de las consultas de similitud (Top-K).

\subsubsection{Métricas de coste computacional}
% Anotación: Definir las métricas de tiempo (latencia por imagen, tiempo total
% de entrenamiento). Mencionar el consumo de recursos (CPU/GPU/Memoria) si se
% han monitorizado.

\subsection{Diseño experimental y procedimiento de pruebas}
\label{subsec:diseno_experimental}

\subsubsection{Entorno de ejecución}
% Anotación: Especificar el hardware utilizado (procesador, tarjeta gráfica, RAM).
% Listar versiones críticas de software (PyTorch, Scikit-learn, PostgreSQL).

\subsubsection{Configuración de escenarios (Whitelists)}
% Anotación: Explicar el funcionamiento de los archivos JSON de "whitelist"
% para orquestar los experimentos.

\subsubsection{Protocolo de pruebas}
% Anotación: Describir los pasos exactos que sigue el pipeline en cada
% ejecución para asegurar condiciones idénticas entre pruebas.

